\uuid{jrkP}
\titre{Limite en un point}
\auteur{bodin ; hanani ; mzari}
\organisation{exo7}
\licence{CC BY-NC-SA 4.0}
\datecreate{2020-01-20}
\niveau{L1}
\module{Analyse}
\chapitre{Continuité, limite et étude de fonctions réelles}
\sousChapitre{Limite de fonctions}
\difficulte{2}
\theme{Limites des fonctions réelles, limite en un point}
\exo7qcmid{lille-1.limites.020}

\contenu{
\begin{question}
\qclassification{123.03}
\qtags{motcle=limite en un point}
Soit $f(x)= \frac{\cos x-1}{x^2}$. Quelles sont les assertions vraies ?
\begin{answers}

    \bad{$f$ n'admet pas de limite en $0$.}

    \bad{$\lim_{x\to 0} f(x)=+\infty$}

    \good{$\lim_{x\to 0} f(x)=-\frac{1}{2}$}

    \bad{$\lim_{x\to 0} f(x)=\frac{1}{2}$}

\end{answers}
\begin{explanations}
On a : $\cos x = \cos^2 (\frac{x}{2}) -  \sin^2 (\frac{x}{2})$ et $1= \cos^2 (\frac{x}{2}) + \sin^2 (\frac{x}{2})$, donc $\cos x - 1 = -2 \sin ^2 (\frac{x}{2})$. D'autre part, $\lim_{x\to 0}\frac{\sin x}{x} = 1$. On déduit que   $\lim_{x\to 0} f(x) = \lim_{x\to 0} -\frac{1}{2}  \big(\frac{\sin (\frac{x}{2})}{\frac{x}{2}}\big)^2 = -\frac{1}{2}$.
\end{explanations}
\end{question}
}
